\section{Related Work}

\paragraph{Knowledge distillation for language models.}
Knowledge distillation transfers a teacher's behavior into a smaller
student~\citep{hinton2015distilling}, and sequence-level distillation extends this idea to autoregressive generation by training on teacher trajectories~\citep{kim2016sequence}. It underlies many strong small instructions and
reasoning models trained from teacher-written solutions or
rationales~\citep{mitra2024orca,yue2023mammoth,yu2023metamath,hsieh2023distilling}.
The standard recipe is efficient but off-policy: the student learns from the teacher
prefixes, then must recover from its own prefixes at inference. On-policy distillation addresses this by querying the teacher on student-sampled prefixes:
MiniLLM uses reverse KL~\citep{gu2024minillm}, GKD mixes student and teacher
generations~\citep{agarwal2024gkd}, DistiLLM refines the divergence and replay
design~\citep{ko2024distillm,ko2025distillm2}, speculative distillation interleaves
teacher and student tokens~\citep{xu2024speculativekd}, and recent practice frames
OPD as dense supervision compared with sparse RL~\citep{thinkingmachines2025onpolicy}.
Recent extensions study black-box OPD and on-policy context
distillation~\citep{ye2025blackboxopd,ye2026opcd}.
By analogy, a stronger teacher is not always a better teacher~\citep{xu2024stronger};
there the concern is teacher and data selection for single-turn instruction tuning,
whereas our reliability notion concerns \emph{where} the teacher's per-step conditional
is queried, so we take it as a suggestive parallel rather than direct support. We extend this line from single-turn generation to
multi-turn training from a fixed teacher-trajectory pool, where the key object is not only the teacher target but the prefix distribution on which it is
queried.

\paragraph{RL algorithms for LLM post-training.}
RL post-training optimizes scalar feedback rather than teacher conditionals. RLHF
fits preference rewards, often Bradley--Terry models~\citep{bradley1952rank}, and
optimizes them with PPO~\citep{schulman2017proximal}, as in instruction-following and
helpful-harmless assistants~\citep{ouyang2022training,bai2022training}. Direct
preference objectives simplify this pipeline by replacing online reward optimization
with contrastive or implicit-reward losses such as SLIC-HF~\citep{zhao2023slic},
DPO~\citep{rafailov2023direct}, IPO~\citep{azar2023general}, and
GPO~\citep{tang2024generalized}, with iterative and online variants studied under
KL-regularized preference learning~\citep{xiong2024iterative,dong2024rlhf}. For
reasoning, o1/DeepSeek-R1-style systems increasingly use verifiable or process
rewards~\citep{jaech2024openai,deepseekai2025deepseekr1incentivizingreasoningcapability,wang2023mathshepherd,zhang2024entropyregularizedprocessrewardmodel},
giving rise to scalable methods such as GRPO~\citep{shao2024deepseekmath} and
DAPO~\citep{yu2025dapo}. Critic-free variants revisit
REINFORCE~\citep{williams1992simple}, including ReMax~\citep{li2023remax},
RLOO-style baselines~\citep{ahmadian2024back,kool2019buy}, and
Reinforce-Rej~\citep{xiong2025minimalist}; RAFT and rejection sampling are the
binary-reward limit, where successful trajectories are selected and
imitated~\citep{dong2023raft,liu2023statistical,xiong2025minimalist}. ReOPD is
orthogonal to this optimizer family: it keeps dense teacher supervision and asks
which replayed prefixes make that supervision reliable.

\paragraph{Reasoning and agentic LLMs.}
RL-style post-training also extends to agents that interact with tools and
environments~\citep{yao2022react,schick2023toolformer,nakano2021webgpt}, including
search-augmented agents~\citep{jin2025search} and tool-integrated mathematical
reasoning~\citep{gou2023tora,mint2024a}. Other work improves reasoning through
self-hinting~\citep{liao2026selfhinting}, budget-aware elastic
reasoning~\citep{xu2025elasticreasoning}, online experiential
learning~\citep{ye2026onlineexperiential}, or distills teacher-generated
reasoning-and-acting traces into smaller agents~\citep{chen2023fireact,chen2024agentflan}.
These works motivate our setting but optimize different signals: scalar rewards,
advantages, filtered rollouts, or fixed teacher trajectories. We retain the dense
conditional target of distillation while avoiding live environment interaction, so the central design variable becomes which replayed prefixes should carry training
mass.

\paragraph{Self-training and rejection sampling.}
Self-training methods repeatedly turn model samples into new training data.
Expert-iteration methods imitate successful search or sampling
solutions~\citep{anthony2017thinking}, STaR bootstraps from self-generated
rationales~\citep{zelikman2022star}, and ReST-style methods alternate generation,
scalar-feedback filtering, and retraining~\citep{gulcehre2023reinforced,singh2023beyond}.
In LLM post-training, rejection sampling and RAFT instantiate the same principle by
imitating reward-ranked trajectories~\citep{dong2023raft,liu2023statistical,xiong2025minimalist},
with refinements for variance, curricula, and agentic
self-improvement~\citep{yao2025optimizing,zhao2024automatic,aksitov2023rest};
reflective methods revise trajectories from feedback~\citep{shinn2024reflexion}. ReOPD also reuses
offline traces, but the trace is not merely accepted or rejected: it becomes a site
where a teacher's condition is queried, so the effective prefix distribution is the
core variable.

\paragraph{Distribution shift and compounding errors.}
The prefix trap inherits the classic sequential-learning problem: behavior cloning on expert states suffers covariate shift once the learner deviates, which DAgger
addresses by training on the learner's induced state distribution~\citep{ross2011reduction}.
Autoregressive generation has the same exposure-bias structure; scheduled sampling,
sequence-level training, and Professor Forcing reduce the mismatch between
teacher-forced training and free-running inference~\citep{bengio2015scheduled,ranzato2016sequence,lamb2016professor}.
LLM alignment revisits this tension at scale: on-policy, even suboptimal, samples can
help preference fine-tuning~\citep{tajwar2024preference}, but on-policy data is not
uniformly best and depends on the alignment stage~\citep{sun2025onpolicydata}. Recent reasoning distillation similarly identifies a dual exposure bias: teacher-forced
traces mismatch student inference, while fully student-generated contexts can make
teacher intervention unreliable~\citep{wang2026backtracking}. Our formulation makes this duality explicit as a two-sided shift: student relevance pulls prefixes toward
the learner, while teacher reliability pulls them back toward the teacher-supported
region.
